% Pre-results Results, Discussion, Limitations, and Conclusion fragment.
% The surrounding manuscript defines the locked architecture, model, dataset,
% and metric commands.  This fallback keeps the fragment compilable when it is
% inspected on its own.
\providecommand{\tbd}[1]{\textit{[TBD: #1]}}

\section{Results}
\label{sec:results}

\paragraph{Evidence status.}
At this stage, the evidence covers implementation behavior rather than
generation quality.  The local \texttt{src/eviseq\_new} suite runs the
value-anchored \AFMR and grounded-copy graph on model-free fixtures.  It checks
source masking and alignment, cross-attention and cache behavior, cross-entropy
gradient reachability, dense and chunked loss parity, checkpoint round trips,
warm-up and full-training paths, and greedy generation with prediction
resumption.  The source documentation also defines a smoke path that trains,
saves, reloads, and evaluates the graph with tiny randomly initialized
backbones.  These checks show that the specified tensor contracts and runtime
paths execute on controlled inputs.  They do not yield a held-out comparison,
a ROUGE-1/2/L value, a factuality estimate, a \BERTScore value, or a hardware
throughput benchmark.

The four-dataset fine-tuning runs and prediction files are not yet available.
The fixed Perl \Rouge-1.5.5 scorer, \AlignScore evaluator, and \BERTScore
evaluator therefore have no results to report.  Dashes in
Table~\ref{tab:results_plan} mean ``not yet computed''; they do not mean zero,
failure, or imputation.  This draft makes no quality claim.

\begin{table*}[t]
\centering
\small
\caption{Pre-results reporting plan.  Each outcome cell remains blank until
the corresponding scorer-locked run and provenance check are complete.}
\label{tab:results_plan}
\begin{tabular}{@{}p{0.31\textwidth}cccc@{}}
\toprule
Planned analysis & \PubMed & \ArXiv & \BookSum & \GovReport \\
\midrule
RQ1: source support (\AlignScore and paired intervals) & -- & -- & -- & -- \\
RQ2: overlap and semantic quality (\Rouge and \BERTScore) & -- & -- & -- & -- \\
RQ3: matched no-encoder, no-bridge, and no-grounded-copy ablations & -- & -- & -- & -- \\
RQ4: encoder-transfer matrix and quality/factuality diagnostics & -- & -- & -- & -- \\
\bottomrule
\end{tabular}
\end{table*}

\paragraph{Evaluation contract.}
The primary comparison uses fine-tuned decoder-only controls:
\QwenZeroSix, \QwenFour, \QwenEight, \LlamaThreeTwoThree,
\LlamaThreeEight, \NemotronThree, and \NemotronEight.  It also includes the
fine-tuned encoder--decoder reference \TfiveGemma.  On each of \PubMed,
\ArXiv, \BookSum, and \GovReport, the runs share example IDs, fixed
train/validation/test splits, source and target budgets, preprocessing,
training-stage budgets, and the registered decoding procedure.  We report
native model serialization and prompt/chat-template behavior instead of
assuming that they match.  One fixed Perl \Rouge-1.5.5 wrapper supplies the
primary \Rouge-1/2/L F1 values; Python \Rouge remains a diagnostic backend and
is never combined with the paper scores.  Each comparison will include a
per-example paired bootstrap interval.  A macro-average across the four
datasets is a secondary summary.  The prediction manifest will record output
and reference lengths, truncation, empty outputs, repetition, source-copy rate,
latency, peak memory, checkpoint identity, and failed or omitted runs.

\paragraph{RQ1: factual-support signal.}
RQ1 asks whether \EviSeq changes an automatic source-support signal relative to
the fine-tuned decoder-only controls on \PubMed, \ArXiv, \BookSum, and
\GovReport.  Its primary quantity is native \AlignScore consistency, where a
higher value means greater source--summary consistency under the metric.  We
may also show $H=1-C$ for $C=\AlignScore_{\texttt{nli\_sp}}$ as a
direction-transformed, lower-is-better quantity for a hallucination-oriented
view.  This transform is not a calibrated probability.  Claims about typed
entity, number, or relation errors require a separately specified adjudicated
audit.

\emph{Fill-in template after the locked runs:} ``Across \PubMed, \ArXiv,
\BookSum, and \GovReport, \EviSeq obtained \AlignScore consistency of
\tbd{dataset-specific values}, compared with \tbd{named matched decoder-only
control(s)}.  The paired differences were \tbd{differences}, with 95\%
bootstrap intervals of \tbd{intervals}.  The corresponding $H=1-C$ values
were \tbd{values}.  This pattern \tbd{supports/does not support} a change in
the automatic source-support signal under this benchmark; it does not by
itself establish a reduction in factual errors.''

\paragraph{RQ2: summary quality.}
RQ2 tests whether a change in source support is accompanied by acceptable
summary coverage and semantic similarity.  It compares \EviSeq with the same
decoder-only controls and \TfiveGemma on \PubMed, \ArXiv, \BookSum, and
\GovReport.  The scorer-locked outcomes are \Rouge-1/2/L F1 and
\BERTScore-F1, read alongside length, truncation, empty-output, and repetition
diagnostics.

\emph{Fill-in template after the locked runs:} ``Against each matched
decoder-only control and \TfiveGemma, \EviSeq changed \Rouge-1, \Rouge-2, and
\Rouge-L F1 on \PubMed, \ArXiv, \BookSum, and \GovReport by
\tbd{dataset-by-baseline deltas}.  The paired 95\% bootstrap intervals were
\tbd{intervals}, and the corresponding \BERTScore differences were
\tbd{differences}.  Output diagnostics showed \tbd{length, truncation,
empty-output, repetition, and copy findings}.  These results \tbd{support/do
not support} a quality change under the registered protocol.''  Here
\BERTScore measures semantic similarity; it is not a hallucination detector.

\paragraph{RQ3: component contribution.}
RQ3 compares the full \EviSeq graph with separately trained no-encoder,
no-\AFMR-bridge, and no-grounded-copy controls.  The decoder and training
budget stay fixed across these runs on \PubMed, \ArXiv, \BookSum, and
\GovReport.  Disabling a component in a finished checkpoint is only a
diagnostic; the primary ablation comparison trains every row from the
pretrained initialization.  The interpretation remains local to the
registered graph and protocol.  Ablation differences can show which
components contribute to the observed benchmark behavior, but they cannot
establish universal causal necessity.

\emph{Fill-in template after the locked runs:} ``On \PubMed, \ArXiv,
\BookSum, and \GovReport, removing the encoder, \AFMR bridge, and grounded
copy route changed the pre-registered \tbd{primary outcome(s)} relative to
full \EviSeq by \tbd{component-specific deltas}.  The paired intervals were
\tbd{intervals}, and the associated \Rouge, \AlignScore, and \BERTScore
diagnostics were \tbd{values}.  The pattern \tbd{does/does not} indicate that
\tbd{component(s)} contributes to the measured behavior under this registered
setup.''

\paragraph{RQ4: encoder portability.}
RQ4 evaluates the same interface with \PPLXEmbed as the anchor and with
\QwenCausal, \QwenEmbed, and \NemotronEmbed encoder variants on \PubMed,
\ArXiv, \BookSum, and \GovReport.  For every variant, the transfer comparison
will report encoder and decoder identifiers, hidden widths, trainable-parameter
counts, and the same quality, factuality, and prediction diagnostics.
Directionality and capacity belong to the encoder choice, so we will keep them
explicit when interpreting transfer differences rather than attributing those
differences to a single cause.

\emph{Fill-in template after the locked runs:} ``With the decoder and
protocol fixed, the \PPLXEmbed, \QwenCausal, \QwenEmbed, and \NemotronEmbed
variants obtained \tbd{\Rouge-1/2/L values}, \tbd{\AlignScore values}, and
\tbd{\BERTScore values} on \PubMed, \ArXiv, \BookSum, and \GovReport.  Their
paired differences from the \PPLXEmbed anchor were \tbd{differences with
intervals}.  Latency, peak memory, length, repetition, and source-copy
diagnostics were \tbd{diagnostics}.  These results \tbd{support/do not
support} portability across the tested encoder choices, conditional on the
observed directionality and capacity differences.''

\section{Discussion}
\label{sec:discussion}

The current result is an implementation boundary.  The source tree provides
an executable \AFMR interface that retains a full retrieval tensor and an
aligned final-state value tensor, separates
refined retrieval keys from final-state values, and exposes a grounded copy
route trained through the same gold-token cross-entropy objective as the
language-model path.  The tests support those execution properties, and the
documented fixture smoke path provides a bounded way to exercise the complete
runtime.  These implementation checks do not answer whether the interface
improves factuality, overlap quality, component-level behavior, or encoder
portability.  Answers to those questions depend on the held-out predictions
and scorer-locked comparisons in Table~\ref{tab:results_plan}.

An \AlignScore increase, accompanied by a decrease in $H=1-C$ for
$C=\AlignScore_{\texttt{nli\_sp}}$,
would be consistent with greater source--summary consistency under the
selected proxy.  It would remain compatible with alternative explanations
such as shorter outputs, stronger extractive behavior, truncation patterns,
or dataset-specific reference properties.  A manually audited factual-error
subset, if available, would be needed to test whether the metric movement
tracks unsupported entities, numbers, and relations.  Copying a source token
does not establish that the token has the correct role in a generated
relation, so source-copy rate cannot replace that audit.

\Rouge and \BERTScore answer a different question.  A change in \Rouge-1/2/L
F1 or \BERTScore-F1 would describe reference overlap or semantic similarity
under the locked preprocessing and scorer.  It would not establish factual
correctness.  Length ratio, empty and truncated predictions, repetition, and
copy diagnostics are needed to determine whether a metric difference reflects
content selection, output constraints, or a change in source use.  A factuality
change with lower overlap would therefore be reported as a possible trade-off
until the paired quality evidence is examined.

The RQ3 ablations provide component attribution only within the tested \EviSeq graph.
Removing a route changes the information path and may also change optimization
dynamics or the number of trainable parameters.  Separate train-from-
pretrained runs, matched budgets, and validation-based development reduce
these concerns but do not turn an ablation into a universal necessity claim.
The narrow interpretation is whether each route contributes to the measured
outcomes under the registered conditions.

The RQ4 matrix tests whether the interface remains usable when the encoder
family, directionality, or capacity changes.  A result that transfers to one
variant would support portability to that tested choice.  It would not isolate
directionality from pretraining data, tokenizer behavior, hidden width, or
parameter count.  These factors should remain visible in the report so that a
transfer effect is not attributed to the \AFMR interface alone.

Several negative-space checks govern the eventual interpretation.  A quality
table without a corresponding prediction manifest would leave checkpoint and
tokenizer provenance unresolved.  An \AlignScore value without its direction
and paired uncertainty would not answer RQ1.  A single aggregate across the
four datasets would conceal a missing dataset-specific result.  A reported
ablation without a matched budget would not identify a component effect.  A
transfer claim without encoder identifiers and capacity diagnostics would not
establish portability.  The final paper should fill each of these evidence
links before replacing the templates above with declarative results.

\section{Limitations}
\label{sec:limitations}

First, the decisive quality evidence is absent at this stage.  No
four-dataset predictions or \Rouge-1.5.5, \AlignScore, or \BERTScore measurements
are available for the proposed \EviSeq configuration.  The implementation
tests cannot support a performance, factuality, or superiority claim.

Second, the present source tree does not include the final BookSum and
GovReport preparation manifests or snapshots.  Their source and target field
choices, split identities, and preprocessing records must be supplied before
the corresponding runs can be treated as reproducible evidence.  Until then,
Table~\ref{tab:results_plan} is a plan rather than a completed four-dataset
evaluation.

Third, the outcome measures have limited scope.  \Rouge measures reference
overlap, and \BERTScore measures semantic similarity; neither is a factuality
test.  \AlignScore is a factuality proxy, and $1-\mathrm{AlignScore}$ is a
direction transform rather than a calibrated hallucination probability.  A
manual factuality audit and multiple training seeds, if added, would address
uncertainty that paired test-example intervals and a metric proxy leave open.

Fourth, source limits and conservative offset alignment can remove or reduce
copy opportunities at truncation boundaries.  Full-memory operation therefore
means that all valid positions within the encoded prefix are retained; it does
not guarantee that information outside the prefix is available or that the
decoder uses every retained position correctly.

Fifth, encoder transfer is bounded by the selected families and by the
confounding of directionality, pretraining, tokenization, width, and
capacity.  Latency and peak-memory costs also require measurement on the
intended hardware.  Tiny local fixtures establish runtime contracts, not
deployment efficiency.

\section{Conclusion}
\label{sec:conclusion}

This paper defines \EviSeq as a source-grounded full-memory interface for
composing a pretrained encoder with a causal decoder.  Its \AFMR bridge keeps
the encoded source available, adapts retrieval keys separately from
final-state values, and supplies a shared source prior; the decoder adds a
grounded copy route while training the complete graph with gold-token
cross-entropy.  These are testable design choices, not a measured quality
claim.

The current evidence supports implementation consistency on controlled
fixtures.  It does not establish an \AlignScore improvement, a \Rouge or
\BERTScore change, component necessity, or encoder portability.  The next
evidence required is the matched, four-dataset evaluation with scorer-locked
predictions, paired uncertainty, ablations, encoder-transfer runs, and the
registered diagnostics.  Until those runs are complete, \EviSeq should be
described as a reproducible interface and evaluation contract whose empirical
contribution remains to be determined.
